% GENERATED FILE. Edit canonical lesson frontmatter, then run npm run generate:books.
% canonical-entries: 55

\chapter*{Glossary}
\addcontentsline{toc}{chapter}{Glossary}
\markboth{Glossary}{Glossary}

This glossary collects every word and phrase taught in the book. Pronunciation appears when it differs from the written form; chapter numbers show where each entry is introduced.

\noindent\begin{minipage}[t]{\linewidth}
\raggedright
\textbf{\pa{ਅੱਖ}}\enspace\emph{akkh}\par
\small eye — a doubled consonant that survived whole, so it never had a tone to lose\par
\footnotesize Introduced in Chapter 12.
\end{minipage}\par\medskip

\noindent\begin{minipage}[t]{\linewidth}
\raggedright
\textbf{\pa{ਆਉਣਾ}}\enspace\emph{āuṇā}\par
\small to come — “go” with “toward” glued on the front, and the “go” half is English come\par
\footnotesize Introduced in Chapter 7.
\end{minipage}\par\medskip

\noindent\begin{minipage}[t]{\linewidth}
\raggedright
\textbf{\pa{ਭੈਣ}}\enspace\emph{bhaiṇ}\par
\small sister — the same low tone as brother, but not the same family of words\par
\footnotesize Introduced in Chapter 11.
\end{minipage}\par\medskip

\noindent\begin{minipage}[t]{\linewidth}
\raggedright
\textbf{\pa{ਭਰਾ}}\enspace\emph{bharā}\par
\small brother — spelled with a breathy bh, said with a plain p and a low tone\par
\footnotesize Introduced in Chapter 11.
\end{minipage}\par\medskip

\noindent\begin{minipage}[t]{\linewidth}
\raggedright
\textbf{\pa{ਚਾਹ}}\enspace\emph{cāh}\par
\small tea — a word that took the overland road out of China, and a level high tone hiding inside its own spelling\par
\footnotesize Introduced in Chapter 10.
\end{minipage}\par\medskip

\noindent\begin{minipage}[t]{\linewidth}
\raggedright
\textbf{\pa{ਧੰਨਵਾਦ}}\enspace\emph{dhannavād}\par
\small thank you — the Sanskritic choice\par
\footnotesize Introduced in Chapter 1.
\end{minipage}\par\medskip

\noindent\begin{minipage}[t]{\linewidth}
\raggedright
\textbf{\pa{ਦਿਲ}}\enspace\emph{dil}\par
\small heart — a Persian loan, not the Sanskrit word Punjabi's other body parts have all used\par
\footnotesize Introduced in Chapter 13.
\end{minipage}\par\medskip

\noindent\begin{minipage}[t]{\linewidth}
\raggedright
\textbf{\pa{ਦੋਸਤ}}\enspace\emph{dost}\par
\small friend — a Persian word that literally means \textquotedblleft{}the one you chose\textquotedblright{}\par
\footnotesize Introduced in Chapter 11.
\end{minipage}\par\medskip

\noindent\begin{minipage}[t]{\linewidth}
\raggedright
\textbf{\pa{ਦੁੱਧ}}\enspace\emph{dudh}\par
\small milk — \textquotedblleft{}that which has been milked,\textquotedblright{} said flat, with no tone at all\par
\footnotesize Introduced in Chapter 10.
\end{minipage}\par\medskip

\noindent\begin{minipage}[t]{\linewidth}
\raggedright
\textbf{\pa{ਹਾਂ} / \pa{ਨਹੀਂ}}\enspace\emph{hā̃ / nahī̃}\par
\small yes / no — heard and understood before they are decoded\par
\footnotesize Introduced in Chapter 1.
\end{minipage}\par\medskip

\noindent\begin{minipage}[t]{\linewidth}
\raggedright
\textbf{\pa{ਹੈ}}\enspace\emph{hai}\par
\small is\par
\footnotesize Introduced in Chapter 2.
\end{minipage}\par\medskip

\noindent\begin{minipage}[t]{\linewidth}
\raggedright
\textbf{\pa{ਹੋਣਾ}}\enspace\emph{hoṇā}\par
\small to be — the infinitive that sounds nothing like \pa{ਹੈ}, because two ancient verbs are braided inside it\par
\footnotesize Introduced in Chapter 7.
\end{minipage}\par\medskip

\noindent\begin{minipage}[t]{\linewidth}
\raggedright
\textbf{\pa{ਇੱਕ} \pa{ਦੋ} \pa{ਤਿੰਨ} \pa{ਚਾਰ} \pa{ਪੰਜ}}\enspace\emph{ikk do tinn chār panj}\par
\small one to five — meeting the panj of Punjabi as an actual number\par
\footnotesize Introduced in Chapter 6.
\end{minipage}\par\medskip

\noindent\begin{minipage}[t]{\linewidth}
\raggedright
\textbf{\pa{ਜਾਣਾ}}\enspace\emph{jāṇā}\par
\small to go — and the one sound law that turns Sanskrit words into Punjabi ones\par
\footnotesize Introduced in Chapter 7.
\end{minipage}\par\medskip

\noindent\begin{minipage}[t]{\linewidth}
\raggedright
\textbf{\pa{ਜਾਣਨਾ}}\enspace\emph{jāṇnā}\par
\small to know — one letter away from “to go,” and the same word as English know\par
\footnotesize Introduced in Chapter 7.
\end{minipage}\par\medskip

\noindent\begin{minipage}[t]{\linewidth}
\raggedright
\textbf{\pa{ਕੰਨ}}\enspace\emph{kann}\par
\small ear — a word that looks like it hides \textquotedblleft{}to do\textquotedblright{} inside it, and does not\par
\footnotesize Introduced in Chapter 12.
\end{minipage}\par\medskip

\noindent\begin{minipage}[t]{\linewidth}
\raggedright
\textbf{\pa{ਖਾਣਾ}}\enspace\emph{khāṇā}\par
\small to eat — and the letter next door that stopped being a consonant and became a tone\par
\footnotesize Introduced in Chapter 7.
\end{minipage}\par\medskip

\noindent\begin{minipage}[t]{\linewidth}
\raggedright
\textbf{\pa{ਖ਼ੁਸ਼ੀ} \pa{ਹੋਈ}}\enspace\emph{khushī hoī}\par
\small pleased to meet you — literally happiness happened\par
\footnotesize Introduced in Chapter 2.
\end{minipage}\par\medskip

\noindent\begin{minipage}[t]{\linewidth}
\raggedright
\textbf{\pa{ਕੀ}}\enspace\emph{kī}\par
\small what\par
\footnotesize Introduced in Chapter 2.
\end{minipage}\par\medskip

\noindent\begin{minipage}[t]{\linewidth}
\raggedright
\textbf{\pa{ਲੈਣਾ}}\enspace\emph{laiṇā}\par
\small to take — a verb that lost a whole consonant on the way here\par
\footnotesize Introduced in Chapter 9.
\end{minipage}\par\medskip

\noindent\begin{minipage}[t]{\linewidth}
\raggedright
\textbf{\pa{ਲਿਖਣਾ}}\enspace\emph{likhṇā}\par
\small to write — from a root meaning to scratch, which is how three unrelated languages named it\par
\footnotesize Introduced in Chapter 8.
\end{minipage}\par\medskip

\noindent\begin{minipage}[t]{\linewidth}
\raggedright
\textbf{\pa{ਮਦਦ} \pa{ਕਰਨਾ}}\enspace\emph{madad karnā}\par
\small to help — literally “to do help,” and the help half is Arabic\par
\footnotesize Introduced in Chapter 9.
\end{minipage}\par\medskip

\noindent\begin{minipage}[t]{\linewidth}
\raggedright
\textbf{\pa{ਮੇਰਾ}}\enspace\emph{merā}\par
\small my\par
\footnotesize Introduced in Chapter 2.
\end{minipage}\par\medskip

\noindent\begin{minipage}[t]{\linewidth}
\raggedright
\textbf{\pa{ਮੇਰਾ} \pa{ਨਾਂ} … \pa{ਹੈ}}\enspace\emph{merā nāṁ … hai}\par
\small my name is…\par
\footnotesize Introduced in Chapter 2.
\end{minipage}\par\medskip

\noindent\begin{minipage}[t]{\linewidth}
\raggedright
\textbf{\pa{ਮੂੰਹ}}\enspace\emph{mūnh}\par
\small mouth, face — the second word this book has found with a level high tone\par
\footnotesize Introduced in Chapter 12.
\end{minipage}\par\medskip

\noindent\begin{minipage}[t]{\linewidth}
\raggedright
\textbf{\pa{ਨੱਕ}}\enspace\emph{nakk}\par
\small nose — the plainest possible cousin of the English word for the very same part\par
\footnotesize Introduced in Chapter 12.
\end{minipage}\par\medskip

\noindent\begin{minipage}[t]{\linewidth}
\raggedright
\textbf{\pa{ਨਾਂ}}\enspace\emph{nāṁ}\par
\small name\par
\footnotesize Introduced in Chapter 2.
\end{minipage}\par\medskip

\noindent\begin{minipage}[t]{\linewidth}
\raggedright
\textbf{\pa{ਨਮਸਤੇ}}\enspace\emph{namaste}\par
\small hello — the broader pan-Indian greeting\par
\footnotesize Introduced in Chapter 1.
\end{minipage}\par\medskip

\noindent\begin{minipage}[t]{\linewidth}
\raggedright
\textbf{\pa{ਪਾਣੀ}}\enspace\emph{pāṇī}\par
\small water — Punjabi's first noun for something you would actually ask for, and the phrase that asks for it politely\par
\footnotesize Introduced in Chapter 10.
\end{minipage}\par\medskip

\noindent\begin{minipage}[t]{\linewidth}
\raggedright
\textbf{\pa{ਪੜ੍ਹਨਾ}}\enspace\emph{paṛhnā}\par
\small to read — and in the language it came from, reading meant saying it out loud\par
\footnotesize Introduced in Chapter 8.
\end{minipage}\par\medskip

\noindent\begin{minipage}[t]{\linewidth}
\raggedright
\textbf{\pa{ਪਰਿਵਾਰ}}\enspace\emph{parivār}\par
\small family — literally \textquotedblleft{}what surrounds you,\textquotedblright{} built from two Sanskrit pieces you can pull apart\par
\footnotesize Introduced in Chapter 11.
\end{minipage}\par\medskip

\noindent\begin{minipage}[t]{\linewidth}
\raggedright
\textbf{\pa{ਪਸੰਦ}}\enspace\emph{pasand}\par
\small liked, pleasing — a Persian noun, not a verb, in the one sentence that will not let you be its subject\par
\footnotesize Introduced in Chapter 9.
\end{minipage}\par\medskip

\noindent\begin{minipage}[t]{\linewidth}
\raggedright
\textbf{\pa{ਪੁੱਛਣਾ}}\enspace\emph{puchhṇā}\par
\small to ask — cousin to the Persian word for asking, and neither borrowed it from the other\par
\footnotesize Introduced in Chapter 9.
\end{minipage}\par\medskip

\noindent\begin{minipage}[t]{\linewidth}
\raggedright
\textbf{\pa{ਰੋਟੀ}}\enspace\emph{roṭī}\par
\small bread — the fourth request, and the retroflex letter this book only mentioned in passing until now\par
\footnotesize Introduced in Chapter 10.
\end{minipage}\par\medskip

\noindent\begin{minipage}[t]{\linewidth}
\raggedright
\textbf{\pa{ਸਮਝਣਾ}}\enspace\emph{samajhṇā}\par
\small to understand — from a root meaning “to wake up,” and the word carries a tone\par
\footnotesize Introduced in Chapter 8.
\end{minipage}\par\medskip

\noindent\begin{minipage}[t]{\linewidth}
\raggedright
\textbf{\pa{ਸਤਿ} \pa{ਸ੍ਰੀ} \pa{ਅਕਾਲ}}\enspace\emph{sat srī akāl}\par
\small hello or goodbye — the Sikh greeting\par
\footnotesize Introduced in Chapter 1.
\end{minipage}\par\medskip

\noindent\begin{minipage}[t]{\linewidth}
\raggedright
\textbf{\pa{ਸ਼ੁਕਰੀਆ}}\enspace\emph{shukrīā}\par
\small thank you — the Perso-Arabic choice\par
\footnotesize Introduced in Chapter 1.
\end{minipage}\par\medskip

\noindent\begin{minipage}[t]{\linewidth}
\raggedright
\textbf{\pa{ਸਿਰ}}\enspace\emph{sir}\par
\small head — a cousin of English horn, not of English head, and an unborrowed twin of Persian sar\par
\footnotesize Introduced in Chapter 13.
\end{minipage}\par\medskip

\noindent\begin{minipage}[t]{\linewidth}
\raggedright
\textbf{\pa{ਸੋਚਣਾ}}\enspace\emph{sochṇā}\par
\small to think — a verb that started out meaning to burn, and then to grieve\par
\footnotesize Introduced in Chapter 8.
\end{minipage}\par\medskip

\noindent\begin{minipage}[t]{\linewidth}
\raggedright
\textbf{\pa{ਤੂੰ} / \pa{ਤੁਸੀਂ}}\enspace\emph{tū̃ / tusī̃}\par
\small you — familiar / respectful\par
\footnotesize Introduced in Chapter 2.
\end{minipage}\par\medskip

\noindent\begin{minipage}[t]{\linewidth}
\raggedright
\textbf{\pa{ਤੁਹਾਡਾ} \pa{ਨਾਂ} \pa{ਕੀ} \pa{ਹੈ}?}\enspace\emph{tuhāḍā nāṁ kī hai?}\par
\small what is your name?\par
\footnotesize Introduced in Chapter 2.
\end{minipage}\par\medskip

\noindent\begin{minipage}[t]{\linewidth}
\raggedright
\textbf{\pa{ਵੇਖਣਾ}}\enspace\emph{vekhṇā}\par
\small to see — built on the old word for the eye itself, the root behind eye, optic, and window\par
\footnotesize Introduced in Chapter 7.
\end{minipage}\par\medskip

\noindent\begin{minipage}[t]{\linewidth}
\raggedright
\textbf{\pa{ਕੋਈ} \pa{ਗੱਲ} \pa{ਨਹੀਂ}}\par
\small it's nothing / no problem / you're welcome\par
\footnotesize Introduced in Chapter 3.
\end{minipage}\par\medskip

\noindent\begin{minipage}[t]{\linewidth}
\raggedright
\textbf{\pa{ਕੰਮ} \pa{ਕਰਨਾ}}\par
\small to work (lit. \textquotedblleft{}work-do\textquotedblright{})\par
\footnotesize Introduced in Chapter 5.
\end{minipage}\par\medskip

\noindent\begin{minipage}[t]{\linewidth}
\raggedright
\textbf{\pa{ਕਿਵੇਂ}}\par
\small how\par
\footnotesize Introduced in Chapter 3.
\end{minipage}\par\medskip

\noindent\begin{minipage}[t]{\linewidth}
\raggedright
\textbf{\pa{ਠੀਕ}}\par
\small fine, well — and the reply \textquotedblleft{}\pa{ਮੈਂ} \pa{ਠੀਕ} \pa{ਹਾਂ}\textquotedblright{}\par
\footnotesize Introduced in Chapter 3.
\end{minipage}\par\medskip

\noindent\begin{minipage}[t]{\linewidth}
\raggedright
\textbf{\pa{ਤੁਸੀਂ} \pa{ਕਿਵੇਂ} \pa{ਹੋ}?}\par
\small how are you? (respectful)\par
\footnotesize Introduced in Chapter 3.
\end{minipage}\par\medskip

\noindent\begin{minipage}[t]{\linewidth}
\raggedright
\textbf{\pa{ਫਿਰ}}\par
\small again, then\par
\footnotesize Introduced in Chapter 4.
\end{minipage}\par\medskip

\noindent\begin{minipage}[t]{\linewidth}
\raggedright
\textbf{\pa{ਫਿਰ} \pa{ਮਿਲਾਂਗੇ}}\par
\small we'll meet again (goodbye)\par
\footnotesize Introduced in Chapter 4.
\end{minipage}\par\medskip

\noindent\begin{minipage}[t]{\linewidth}
\raggedright
\textbf{\pa{ਬੋਲਣਾ}}\par
\small to speak\par
\footnotesize Introduced in Chapter 5.
\end{minipage}\par\medskip

\noindent\begin{minipage}[t]{\linewidth}
\raggedright
\textbf{\pa{ਮੈਂ}}\par
\small I\par
\footnotesize Introduced in Chapter 3.
\end{minipage}\par\medskip

\noindent\begin{minipage}[t]{\linewidth}
\raggedright
\textbf{\pa{ਮੈਂ} \pa{ਪੰਜਾਬੀ} \pa{ਬੋਲਦਾ} \pa{ਹਾਂ}}\par
\small I speak Punjabi (gendered)\par
\footnotesize Introduced in Chapter 5.
\end{minipage}\par\medskip

\noindent\begin{minipage}[t]{\linewidth}
\raggedright
\textbf{\pa{ਮਿਲਾਂਗੇ}}\par
\small (we) will meet\par
\footnotesize Introduced in Chapter 4.
\end{minipage}\par\medskip

\noindent\begin{minipage}[t]{\linewidth}
\raggedright
\textbf{\pa{ਰੱਬ} \pa{ਰਾਖਾ}}\par
\small goodbye (lit. \textquotedblleft{}God {[}be your{]} keeper\textquotedblright{})\par
\footnotesize Introduced in Chapter 4.
\end{minipage}\par\medskip

\noindent\begin{minipage}[t]{\linewidth}
\raggedright
\textbf{\pa{ਰਹਿਣਾ}}\par
\small to live, to stay\par
\footnotesize Introduced in Chapter 5.
\end{minipage}\par\medskip
